% Sauer–Shelah 引理证明中的“删除末坐标、按前缀分组”（第5章5.2.4）。
% 模式族 A = {0,1}^3 \ {101}，共 7 个模式，VC 维（可打散的最大坐标数）为 2。
% 删除第 3 坐标后得到前缀：00 接 {0,1}，01 接 {0,1}，10 只接 0，11 接 {0,1}。
% U = {00,01,10,11}（|U|=4），W = {00,01,11}（可接两种末位的前缀，|W|=3），|A| = |U| + |W| = 7。
% 单元格 5x5 mm；被删除的第 3 坐标用黄色浅底标出。
\begingroup
\begin{tikzpicture}[figure picture,x=1mm,y=1mm,
  bit/.style={draw=FigureStroke,line width=.6pt,minimum size=5mm,inner sep=0pt,anchor=south west},
  last/.style={bit,fill=FigureYellowFill},
  pre/.style={bit,fill=FigureBlueFill}]
\node[figure heading] at (7.5,51) {模式族 $\mathcal A$};
\node[figure heading] at (37.5,51) {前缀};
\node[figure heading] at (58,51) {可接的末位};
\node[figure heading] at (78,51) {$U$};
\node[figure heading] at (92,51) {$W$};
% groups: prefix/bits/rows-y
\foreach \p/\q/\ya/\yb/\yc/\tail/\inw in {
  0/0/41/35/38/{0, 1}/1,
  0/1/27/21/24/{0, 1}/1,
  1/0/13/13/13/{只接 0}/0,
  1/1/5/-1/2/{0, 1}/1}{
  \node[bit] at (0,\ya) {$\p$};\node[bit] at (5,\ya) {$\q$};
  \node[last] at (10,\ya) {$0$};
  \ifnum\ya=\yb\else
    \node[bit] at (0,\yb) {$\p$};\node[bit] at (5,\yb) {$\q$};
    \node[last] at (10,\yb) {$1$};
    \draw[FigureMuted] (16.5,\ya+4) -- (17.5,\ya+4) -- (17.5,\yb+1) -- (16.5,\yb+1);
  \fi
  \draw[figure flow] (18.5,\yc+2.5) -- (32,\yc+2.5);
  \node[pre] at (33,\yc) {$\p$};\node[pre] at (38,\yc) {$\q$};
  \node at (58,\yc+2.5) {\tail};
  \node[text=FigureBlue] at (78,\yc+2.5) {$\checkmark$};
  \ifnum\inw=1 \node[text=FigureBlue] at (92,\yc+2.5) {$\checkmark$};
  \else \node[text=FigureMuted] at (92,\yc+2.5) {--};\fi
}
\node[figure note,anchor=west] at (0,-6.5) {缺少 $101$；黄色列为删去的末坐标};
\draw[FigureGrid] (68,-3) -- (100,-3);
\node at (78,-7) {$|U|=4$};\node at (92,-7) {$|W|=3$};
\node[figure note,anchor=east] at (100,-12.5) {$|\mathcal A|=|U|+|W|=7$};
\end{tikzpicture}
\endgroup
